\documentclass[a4paper]{article}

% 文章排版
\usepackage{indentfirst}
\setlength{\parindent}{0em}
\usepackage{autobreak}
\allowdisplaybreaks
\usepackage{parskip}
\usepackage{geometry}
\geometry{top=3cm,bottom=3cm,left=2cm,right=2cm}
\usepackage{anyfontsize}
\usepackage{hyperref}

% 数学宏包
\usepackage{amsmath}
\usepackage{esint}

% 设置数学字体

% 如果使用 unicode-math 包，需要将 \mathbf 命令换成 \symbf 命令
% 且对 \mathscr 和 \mathcal 的粗体直接使用 \mathbfscr 和 \mathbfcal

% Latin Modern 风格设置（默认）
\usepackage{unicode-math}
\unimathsetup{mathrm=sym,mathbf=sym,mathsf=sym,math-style=ISO,bold-style=ISO,sans-style=italic}
\usepackage[scr=rsfs]{mathalpha}

% Times New Roman 风格设置
% \usepackage{fontspec}
% \setmainfont{XITS}
% \usepackage{unicode-math}
% \unimathsetup{mathrm=sym,mathbf=sym,mathsf=sym,math-style=ISO,bold-style=ISO,sans-style=italic}
% \setmathfont{XITS Math}
% \setmathfont[range={\mathcal,\mathbfcal},StylisticSet=1]{XITS Math}

% 第二种 Times New Roman 风格设置
% \usepackage{fontspec}
% \setmainfont{STIX Two Text}
% \usepackage{unicode-math}
% \unimathsetup{mathrm=sym,mathbf=sym,mathsf=sym,math-style=ISO,bold-style=ISO,sans-style=italic}
% \setmathfont{STIX Two Math}
% \setmathfont[range={\mathscr,\mathbfscr},StylisticSet=1]{STIX Two Math}

% 第三种 Times New Roman 风格设置
% 不需要 fontspec 和 unicode-math
% \usepackage{newtxtext}
% \usepackage[slantedGreek,vvarbb]{newtxmath}
% \defaultfontfeatures{}
% \usepackage{bm}
% \renewcommand\mathbf\bm

% Linux Libertine 风格设置
% 不需要 fontspec 和 unicode-math
% \usepackage[libertine]{newtx}
% \usepackage{bm}
% \renewcommand\mathbf\bm

% 第二种 Linux Libertine 风格设置（有一些小的 kerning 问题）
% 不需要 fontspec
% \usepackage{unicode-math}
% \unimathsetup{mathrm=sym,mathbf=sym,mathsf=sym,bold-style=ISO,sans-style=italic}
% \usepackage[math=+ss08]{libertinus-otf}
% \setmathfont[range={\mathcal,\mathbfcal}]{Latin Modern Math}

% Palatino 风格设置
% 不需要 unicode-math
% \usepackage{fontspec}
% \usepackage{newpxtext}
% \usepackage[slantedGreek,vvarbb]{newpxmath}
% \defaultfontfeatures{}
% \renewcommand\hbar\hslash
% \usepackage{bm}
% \renewcommand\mathbf\bm

% 第二种 Palatino 风格设置
% 不需要 fontspec 和 unicode-math
% \usepackage[T1]{fontenc}
% \usepackage[slantedGreek]{mathpazo}
% \renewcommand\mathbf\mathbold
% \usepackage[scr=kp,frak=pxtx]{mathalpha}

% 第三种 Palatino 风格设置（不完全是 Palatino）
% 不需要 fontspec
% \usepackage{unicode-math}
% \unimathsetup{mathrm=sym,mathbf=sym,mathsf=sym,math-style=ISO,bold-style=ISO,sans-style=italic}
% \usepackage{kpfonts-otf}

% Cambria 风格设置
% \usepackage{fontspec}
% \setmainfont{Cambria}[BoldFont=CambriaB,BoldItalicFont=CambriaZ]
% \usepackage{unicode-math}
% \unimathsetup{mathrm=sym,mathbf=sym,mathsf=sym,math-style=ISO,bold-style=ISO,sans-style=italic}
% \setmathfont{Cambria Math}
% \setmathfont[range={\mathscr,\mathbfscr},StylisticSet=1]{STIX Two Math}

% Lucida 风格设置
% \usepackage{fontspec}
% \setmainfont{Lucida Bright OT}
% \setmonofont{Lucida Console DK}
% \usepackage{unicode-math}
% \unimathsetup{mathrm=sym,mathbf=sym,mathsf=sym,math-style=ISO,bold-style=ISO,sans-style=italic}
% \setmathfont[StylisticSet=4]{Lucida Bright Math OT}
% \setmathfont[range={\mathscr,\mathbfscr}]{Lucida Bright Math OT}

% Garamond 风格设置
% \usepackage{fontspec}
% \setmainfont{EB Garamond}
% \usepackage{unicode-math}
% \unimathsetup{mathrm=sym,mathbf=sym,mathsf=sym,math-style=ISO,bold-style=ISO,sans-style=italic}
% \setmathfont[StylisticSet={3,6,9,10}]{Garamond-Math}
% \setmathfont[range={\mathscr,\mathbfscr}]{Garamond-Math}

% XCharter 风格设置
% \usepackage{fontspec}
% \setmainfont{XCharter}
% \usepackage{unicode-math}
% \unimathsetup{mathrm=sym,mathbf=sym,mathsf=sym,math-style=ISO,bold-style=ISO,sans-style=italic}
% \setmathfont{XCharter Math}
% \setmathfont[range={\mathcal,\mathbfcal}]{Latin Modern Math}

% 中文
% \usepackage[fontset=none]{ctex}
% \setCJKmainfont{Adobe Song Std}[BoldFont=Noto Sans CJK SC Medium,ItalicFont=Adobe KaiTi Std]
% \setCJKmonofont{Noto Sans Mono CJK SC}[BoldFont=Noto Sans Mono CJK SC Bold,ItalicFont=Adobe FangSong Std]

% \epsilon 符号重定义
\AtBeginDocument{
  \let\uglyepsilon\epsilon
  \renewcommand\epsilon\varepsilon
}

\title{\textbf{Maths Homework Template}}
\author{Author}
\date{\today}

\begin{document}

    \maketitle

    \textbf{Short Math Guide: \url{http://tug.ctan.org/info/short-math-guide}.}

    \textbf{Short Math Guide (Chinese version): \url{https://github.com/hoganbin/short-math-guide}}.

    \bigskip
    \begin{enumerate}

        \item
        \textparagraph\ Inline equations for simple math: $\mathbf{p} = \hbar\mathbf{k}$, $\epsilon = \hbar\mathbf{\omega}$.

        \textparagraph\ Centered equations for more involved or important equations:
        \begin{equation}
            \mathbf{a} \cdot \mathbf{b} = \sum_{i=1}^n a_i b_i
        \end{equation}
        % 如果想要公式编号重新从 1 开始，可以在该公式前使用命令：
        % \setcounter{equation}{0}
        \begin{equation}
            \frac{\mathrm{d}}{\mathrm{d}x}\mathrm{e}^{a(x)} = a'(x)\mathrm{e}^{a(x)}
        \end{equation}
        \begin{equation}
        \frac{\|\mathsf{A}\mathbf{x}\|_\infty}{\|\mathbf{x}\|_\infty} = \frac{\max\limits_i\left|\sum\limits_{j=1}^n \mathsf{A}_{ij}\mathsf{x}_{j1}\right|}{\max\limits_i|\mathsf{x}_{i1}|}
        \end{equation}
        \textparagraph\ An example of a matrix:
        \begin{equation}
            \Psi =
            \begin{bmatrix}
                \dfrac{1+\sqrt{3}i}{2} & \mathrm{e}^{-i\omega t} \\
                i\mathrm{e}^{-i\omega t} & \dfrac{1-\sqrt{3}i}{2}
            \end{bmatrix}
        \end{equation}
        \textparagraph\ An example of intergals:
        \begin{equation}
        \iiint \nabla \cdot \mathbf{A} \,\mathrm{d}V = \oiint \mathbf{A}\,\mathrm{d}\mathbf{S}
        \end{equation}
        \textparagraph\ An example of roots:
        \begin{equation}
        r = \sqrt[3]{\frac{3}{4\pi N}}
        \end{equation}

        \item
        \begin{align*}
            I(x)
            &= \int_{-\tfrac{b}{2}}^{\tfrac{b}{2}}\mathrm{d}I(x) \\
            &= 4\alpha\frac{I_0}{b}\int_{-\tfrac{b}{2}}^{\tfrac{b}{2}}\cos^2\left(\frac{\pi d}{\lambda D}\left(x+\frac{D}{R}\xi\right)\right)\mathrm{d}\xi \\
            &= 4\alpha I_0\left(1+\frac{1}{u}\sin u\cos\left(\frac{2\pi d}{\lambda D}x\right)\right),\quad u = \dfrac{\pi db}{\lambda R}
        \end{align*}

        \item
        \[
        Y_n(y) = \left\{
        \begin{matrix}
            &0, &n = 2k \\
            &\dfrac{8a^2}{\pi^3 n^3}\left(1 - \dfrac{\cosh\dfrac{n\pi y}{a}}{\cosh\dfrac{n\pi b}{2a}}\right), &n = 2k + 1
        \end{matrix}
        \right.
        \]

        \item
        \[
        \begin{matrix}
            &\mathcal{ABCD} &\mathscr{ABCD} &\mathfrak{ABCD} &\mathbb{ABCD} \\
            &\mathsf{ABCD} &\mathbf{ABCD} &\mathrm{ABCD} &\mathit{ABCD}
        \end{matrix}
        \]

        \item
        \[
            u(x,t) = \sum_{n=0}^\infty X_n(x)Y_n(y) = \frac{8a^2}{\pi^3}\sum_{k=0}^\infty \frac{1}{(2k+1)^3}\left(1 - \frac{\cosh\dfrac{2k+1}{a}\pi y}{\cosh\dfrac{2k+1}{2a}\pi b}\right)\sin\frac{2k+1}{a}\pi x
        \]

        \item
        \begin{align*}
            \mathscr{L}\left\lbrace\frac{1-\cos\omega t}{t}\right\rbrace
            &=\int_{p}^{\infty}\left(\frac{1}{p}-\frac{1}{2(p+\mathrm{i}\omega)}-\frac{1}{2(p-\mathrm{i}\omega)}\right)\mathrm{d}t \\
            &=\ln p-\frac{1}{2}\ln(p+\mathrm{i}\omega)-\frac{1}{2}\ln(p-\mathrm{i}\omega)\bigg|_{p}^{\infty} \\
            &=\frac{1}{2}\ln\frac{p^2+\omega^2}{p^2}
        \end{align*}

        \item
        \begin{align*}
            [q,p]_{Q,P}
            &= \frac{\partial q}{\partial Q}\cdot\frac{\partial p}{\partial P} - \frac{\partial q}{\partial P}\cdot\frac{\partial p}{\partial Q} \\
            &= 1 + \frac{P^2\mathrm{e}^{2Q}}{1-P^2\mathrm{e}^{2Q}} - \frac{P^2\mathrm{e}^{2Q}}{1-P^2\mathrm{e}^{2Q}} \\
            &= 1
        \end{align*}

        \item
        \[
        \frac{4\pi R_{\odot}^2 \cdot \sigma T_{\odot}^4}{4\pi d^2} = 1352.83\,\mathrm{W/m^2}
        \]
        \[
        T_{\odot} = 5763.43\,\mathrm{K}
        \]
        \[
        \lambda_{\odot} = \frac{b}{T_{\odot}} = 5027.86\,\text{\AA}
        \]

    \end{enumerate}

\end{document}
